\documentclass[10pt,a4paper]{amsart}
\usepackage{amsmath,amssymb,amsthm,graphicx}
% Retained minimal preamble for the Stacks Project chapter more-algebra.tex (More on Algebra).
% The chapter file pinned-original-source.tex carries only \input{preamble} on its first line, so no
% preamble text lives inside the chapter itself.  The Stacks Project uses one single shared preamble for
% every chapter; the snapshot stages it identically as sources/stacks-chow/preamble.tex and
% sources/stacks-etale-cohomology/preamble.tex.  The block below is copied mechanically, line for line,
% from sources/stacks-chow/preamble.tex lines 14-18 (the Xy-pic setup) and 154-217 (the native theorem
% declarations together with the mathematical macros).  Nothing was added, removed or redefined; the
% wrapper class and the amsmath/amssymb/amsthm/graphicx base packages are supplied by the assembler.
% Chapter: more-algebra.tex -- More on Algebra.
% For commutative diagrams we use Xy-pic
\usepackage[all]{xy}

% We use 2cell for 2-commutative diagrams.
\xyoption{2cell}

% Theorem environments.
%
\theoremstyle{plain}
\newtheorem{theorem}[subsection]{Theorem}
\newtheorem{proposition}[subsection]{Proposition}
\newtheorem{lemma}[subsection]{Lemma}

\theoremstyle{definition}
\newtheorem{definition}[subsection]{Definition}
\newtheorem{example}[subsection]{Example}
\newtheorem{exercise}[subsection]{Exercise}
\newtheorem{situation}[subsection]{Situation}

\theoremstyle{remark}
\newtheorem{remark}[subsection]{Remark}
\newtheorem{remarks}[subsection]{Remarks}

\numberwithin{equation}{subsection}

% Macros
%
\def\lim{\mathop{\mathrm{lim}}\nolimits}
\def\colim{\mathop{\mathrm{colim}}\nolimits}
\def\Spec{\mathop{\mathrm{Spec}}}
\def\Hom{\mathop{\mathrm{Hom}}\nolimits}
\def\Ext{\mathop{\mathrm{Ext}}\nolimits}
\def\SheafHom{\mathop{\mathcal{H}\!\mathit{om}}\nolimits}
\def\SheafExt{\mathop{\mathcal{E}\!\mathit{xt}}\nolimits}
\def\Sch{\mathit{Sch}}
\def\Mor{\mathop{\mathrm{Mor}}\nolimits}
\def\Ob{\mathop{\mathrm{Ob}}\nolimits}
\def\Sh{\mathop{\mathit{Sh}}\nolimits}
\def\NL{\mathop{N\!L}\nolimits}
\def\CH{\mathop{\mathrm{CH}}\nolimits}
\def\proetale{{pro\text{-}\acute{e}tale}}
\def\etale{{\acute{e}tale}}
\def\QCoh{\mathit{QCoh}}
\def\Ker{\mathop{\mathrm{Ker}}}
\def\Im{\mathop{\mathrm{Im}}}
\def\Coker{\mathop{\mathrm{Coker}}}
\def\Coim{\mathop{\mathrm{Coim}}}

% Boxtimes
%
\DeclareMathSymbol{\boxtimes}{\mathbin}{AMSa}{"02}

%
% Macros for moduli stacks/spaces
%
\def\QCohstack{\mathcal{QC}\!\mathit{oh}}
\def\Cohstack{\mathcal{C}\!\mathit{oh}}
\def\Spacesstack{\mathcal{S}\!\mathit{paces}}
\def\Quotfunctor{\mathrm{Quot}}
\def\Hilbfunctor{\mathrm{Hilb}}
\def\Curvesstack{\mathcal{C}\!\mathit{urves}}
\def\Polarizedstack{\mathcal{P}\!\mathit{olarized}}
\def\Complexesstack{\mathcal{C}\!\mathit{omplexes}}
% \Pic is the operator that assigns to X its picard group, usage \Pic(X)
% \Picardstack_{X/B} denotes the Picard stack of X over B
% \Picardfunctor_{X/B} denotes the Picard functor of X over B
\def\Pic{\mathop{\mathrm{Pic}}\nolimits}
\def\Picardstack{\mathcal{P}\!\mathit{ic}}
\def\Picardfunctor{\mathrm{Pic}}
\def\Deformationcategory{\mathcal{D}\!\mathit{ef}}

% --- appended 2026-10-12 (Chief): the project own definitions of the reference/slogan/history
% display environments, copied mechanically from the shared Stacks preamble lines 8-12
% (\usepackage{verbatim} plus the three \newenvironment lines, which the chapter preamble keeps
% verbatim from the project preamble). Nothing added, removed or redefined: these environments
% carry no mathematical content and the project itself typesets them as comments. Without these
% three lines a counted statement whose own text uses them fails with Environment slogan undefined.
\usepackage{verbatim}
\newenvironment{reference}{\comment}{\endcomment}
\newenvironment{slogan}{\comment}{\endcomment}
\newenvironment{history}{\comment}{\endcomment}
\usepackage{cleveref}
\crefname{originallemma}{original lemma}{original lemmas}
\Crefname{originallemma}{Original lemma}{Original lemmas}
\makeatletter
\newlabel{derived-lemma-morphisms-from-projective-complex}{{064B}{source}{}{sourcepointer-1}{}}
\newlabel{derived-lemma-morphisms-from-projective-complex@cref}{{[originallemma][0][]064B}{[1][0][]source}{}{}{}}
\makeatother
\crefname{originallemma}{original lemma}{original lemmas}
\Crefname{originallemma}{Original lemma}{Original lemmas}
\makeatletter
\newlabel{derived-section-ext}{{06XP}{source}{}{sourcepointer-2}{}}
\newlabel{derived-section-ext@cref}{{[originallemma][0][]06XP}{[1][0][]source}{}{}{}}
\makeatother
\crefname{originallemma}{original lemma}{original lemmas}
\Crefname{originallemma}{Original lemma}{Original lemmas}
\makeatletter
\newlabel{item-compatible}{{18533}{source}{}{sourcepointer-3}{}}
\newlabel{item-compatible@cref}{{[originallemma][0][]18533}{[1][0][]source}{}{}{}}
\makeatother
\crefname{originallemma}{original lemma}{original lemmas}
\Crefname{originallemma}{Original lemma}{Original lemmas}
\makeatletter
\newlabel{lemma-RHom-out-of-projective}{{18708}{source}{}{sourcepointer-4}{}}
\newlabel{lemma-RHom-out-of-projective@cref}{{[originallemma][0][]18708}{[1][0][]source}{}{}{}}
\makeatother
\crefname{originallemma}{original lemma}{original lemmas}
\Crefname{originallemma}{Original lemma}{Original lemmas}
\makeatletter
\newlabel{lemma-evaluate-and-more}{{18164}{source}{}{sourcepointer-5}{}}
\newlabel{lemma-evaluate-and-more@cref}{{[originallemma][0][]18164}{[1][0][]source}{}{}{}}
\makeatother
\crefname{originallemma}{original lemma}{original lemmas}
\Crefname{originallemma}{Original lemma}{Original lemmas}
\makeatletter
\newlabel{lemma-internal-hom}{{18682}{source}{}{sourcepointer-6}{}}
\newlabel{lemma-internal-hom@cref}{{[originallemma][0][]18682}{[1][0][]source}{}{}{}}
\makeatother
\crefname{originallemma}{original lemma}{original lemmas}
\Crefname{originallemma}{Original lemma}{Original lemmas}
\makeatletter
\newlabel{lemma-left-dual-complex}{{18462}{source}{}{sourcepointer-7}{}}
\newlabel{lemma-left-dual-complex@cref}{{[originallemma][0][]18462}{[1][0][]source}{}{}{}}
\makeatother
\crefname{originalsection}{original section}{original sections}
\Crefname{originalsection}{Original section}{Original sections}
\makeatletter
\newlabel{section-sign-rules}{{18339}{source}{}{sourcepointer-8}{}}
\newlabel{section-sign-rules@cref}{{[originalsection][0][]18339}{[1][0][]source}{}{}{}}
\makeatother
\title{The dual of a perfect complex is perfect and refines the double dual and the tensor-Hom adjunction}
\author{The Stacks Project Authors}
\date{stacks-project:more-algebra.tex:local-snapshot-20261012; GNU FDL 1.2 or later}
\begin{document}
\maketitle
\setcounter{subsection}{1}
\noindent\textit{Editorial note.} All delivered mathematical text is the author's own, copied from the pinned Stacks Project chapter file more-algebra.tex; no statement or proof step was rewritten, reordered or removed.\par\medskip
\noindent\textit{Editorial note.} The chapter file pinned-original-source.tex carries only \textbackslash{}input\{preamble\} on its first line, so no preamble text lives inside the chapter. The retained file sources/stacks-more-algebra/retained-preamble.tex is a mechanical, line-for-line copy of the project's single shared preamble as staged in the snapshot; nothing was added, removed or redefined and no package was added.\par\medskip
\noindent\textit{Editorial note.} The source URL is the source-level Stacks Project URL, as required when more than one card is produced from one source; the per-claim tag is carried by claim\_key.\par\medskip
\noindent\textit{Editorial note.} The author writes '(this uses Lemma ... to see that E\textasciicircum{}bullet is the left dual of K\textasciicircum{}bullet in the category of complexes)' inside the counted proof span, reproduced verbatim. The nine references inside the counted proof span are recorded as non-load-bearing navigation pointers; they include cross-chapter tags and same-chapter label lines.\par\medskip
\noindent\textit{Allowed standard prerequisites (editorial).}\begin{itemize}
\item Bounded complexes of finite projective modules representing perfect objects. used at proof lines 19189-19190. Is the representation on which the explicit computation is carried out.
\item The internal Hom complex and the evaluation map up to sign. invoked at proof lines 19196-19202. Supplies the canonical map whose bijectivity gives the double dual isomorphism.
\item Left duals in the category of complexes. invoked at proof lines 19220-19222. Is what turns the sign-compatible isomorphism into the tensor-Hom adjunction.
\end{itemize}\medskip
\ifdefined\hypertarget\hypertarget{sourcepointer-1}{}\fi
\noindent\textit{Source reference.} Original lemma 064B: focused-seeds/AGP/sources/stacks-snapshot-20261011/tags.txt, source lines 7860--7860. Cross-chapter navigation pointer into the Stacks Project; the standard fact it supplies is applied inside the counted proof span and the argument resting on it stays there.\par\medskip
\ifdefined\hypertarget\hypertarget{sourcepointer-2}{}\fi
\noindent\textit{Source reference.} Original lemma 06XP: focused-seeds/AGP/sources/stacks-snapshot-20261011/tags.txt, source lines 8858--8858. Cross-chapter navigation pointer into the Stacks Project; the standard fact it supplies is applied inside the counted proof span and the argument resting on it stays there.\par\medskip
\ifdefined\hypertarget\hypertarget{sourcepointer-3}{}\fi
\noindent\textit{Source reference.} Original lemma 18533: sources/stacks-more-algebra/pinned-original-source.tex, source lines 18533--18533. Same-chapter navigation pointer; the statement it names is standard for the section and the argument that consumes it stays inside the counted proof span.\par\medskip
\ifdefined\hypertarget\hypertarget{sourcepointer-4}{}\fi
\noindent\textit{Source reference.} Original lemma 18708: sources/stacks-more-algebra/pinned-original-source.tex, source lines 18708--18708. Same-chapter navigation pointer; the statement it names is standard for the section and the argument that consumes it stays inside the counted proof span.\par\medskip
\ifdefined\hypertarget\hypertarget{sourcepointer-5}{}\fi
\noindent\textit{Source reference.} Original lemma 18164: sources/stacks-more-algebra/pinned-original-source.tex, source lines 18164--18164. Same-chapter navigation pointer; the statement it names is standard for the section and the argument that consumes it stays inside the counted proof span.\par\medskip
\ifdefined\hypertarget\hypertarget{sourcepointer-6}{}\fi
\noindent\textit{Source reference.} Original lemma 18682: sources/stacks-more-algebra/pinned-original-source.tex, source lines 18682--18682. Same-chapter navigation pointer; the statement it names is standard for the section and the argument that consumes it stays inside the counted proof span.\par\medskip
\ifdefined\hypertarget\hypertarget{sourcepointer-7}{}\fi
\noindent\textit{Source reference.} Original lemma 18462: sources/stacks-more-algebra/pinned-original-source.tex, source lines 18462--18462. Same-chapter navigation pointer; the statement it names is standard for the section and the argument that consumes it stays inside the counted proof span.\par\medskip
\ifdefined\hypertarget\hypertarget{sourcepointer-8}{}\fi
\noindent\textit{Source reference.} Original section 18339: sources/stacks-more-algebra/pinned-original-source.tex, source lines 18339--18339. Same-chapter navigation pointer; the statement it names is standard for the section and the argument that consumes it stays inside the counted proof span.\par\medskip
\begin{lemma}
\label{lemma-dual-perfect-complex}
Let $A$ be a ring. Let $K \in D(A)$ be perfect. Then $K^\vee = R\Hom_A(K, A)$
is a perfect complex and $K \cong (K^\vee)^\vee$. There are functorial
isomorphisms
$$
L \otimes_A^\mathbf{L} K^\vee = R\Hom_A(K, L)
\quad\text{and}\quad
H^0(L \otimes_A^\mathbf{L} K^\vee) = \Ext_A^0(K, L)
$$
for $L \in D(A)$.
\end{lemma}

\begin{proof}
We can represent $K$ by a bounded complex $K^\bullet$ of finite projective
$A$-modules. By Lemma \ref{lemma-RHom-out-of-projective} the object $K^\vee$
is represented by the complex $E^\bullet = \Hom^\bullet(K^\bullet, A)$.
Note that $E^n = \Hom_A(K^{-n}, A)$ is finite projective as well.
Thus $E^\bullet$ is a bounded complex of finite projective modules
and we conclude $K^\vee$ is perfect. There is a canonical map
$$
K^\bullet = \text{Tot}(\Hom^\bullet(A, A) \otimes_A K^\bullet)
\longrightarrow
\Hom^\bullet(\Hom^\bullet(K^\bullet, A), A)
$$
which up to sign uses the evaluation map in each degree,
see Lemma \ref{lemma-evaluate-and-more}. (For sign rules see
Section \ref{section-sign-rules}.)
Thus this map defines a canonical isomorphism $(K^\vee)^\vee \cong K$
as the double dual of a finite projective module is itself.

\medskip\noindent
The second equality follows from the first by
Lemma \ref{lemma-internal-hom} and
Derived Categories, Lemma \ref{derived-lemma-morphisms-from-projective-complex}
as well as the definition of Ext groups, see
Derived Categories, Section \ref{derived-section-ext}.
Let $L^\bullet$ be a complex of $A$-modules representing $L$.
By Section \ref{section-sign-rules} item (\ref{item-compatible})
there is a canonical isomorphism
$$
\text{Tot}(L^\bullet \otimes_A E^\bullet)
\longrightarrow
\Hom^\bullet(K^\bullet, L^\bullet)
$$
of complexes of $A$-modules (this uses Lemma \ref{lemma-left-dual-complex}
to see that $E^\bullet$ is the left dual of $K^\bullet$ in the category
of complexes). This proves the first displayed equality
and the proof is complete.
\end{proof}

\end{document}
